% !TeX program = tectonic
\documentclass[11pt,a4paper]{article}
\usepackage{fontspec}
\usepackage[english]{babel}
\babelfont{rm}[Language=Default]{Liberation Serif}
\babelfont[Language=Default]{sf}{Carlito}
\babelfont[Language=Default]{tt}{DejaVu Sans Mono}
\usepackage{amsmath,amssymb,amsthm}
\usepackage{geometry}
\geometry{a4paper, top=2.4cm, bottom=2.4cm, left=2.4cm, right=2.4cm}
\usepackage{booktabs,tabularx,enumitem,xcolor,tcolorbox}
\tcbuselibrary{breakable,skins}
\definecolor{linkblue}{HTML}{1F4E79}
\definecolor{statgreen}{HTML}{2F6B3A}
\definecolor{goldacc}{HTML}{8B7E5A}
\usepackage[colorlinks=true,linkcolor=linkblue,urlcolor=linkblue,unicode=true]{hyperref}
\theoremstyle{plain}
\newtheorem{theorem}{Theorem}
\newtheorem{lemma}[theorem]{Lemma}
\newtheorem{proposition}[theorem]{Proposition}
\newtheorem{corollary}[theorem]{Corollary}
\theoremstyle{definition}
\newtheorem{definition}[theorem]{Definition}
\theoremstyle{remark}
\newtheorem{remark}[theorem]{Remark}
\newtcolorbox{statusbox}[1][]{colback=statgreen!5,colframe=statgreen!75,
  title={\textbf{Summary of results:} #1},fonttitle=\small,boxrule=0.7pt,arc=2pt,breakable}
\newtcolorbox{databox}[1][]{colback=goldacc!6,colframe=goldacc!80,
  title={\textbf{Protocol data:} #1},fonttitle=\small,boxrule=0.7pt,arc=2pt,breakable}
\newtcolorbox{verifybox}[1][]{colback=linkblue!5,colframe=linkblue!80,
  title={\textbf{Verification:} #1},fonttitle=\small,boxrule=0.7pt,arc=2pt,breakable}
\widowpenalty=10000
\clubpenalty=10000
\setlength{\parskip}{0.35em}
\newcommand{\CC}{\mathbb{C}}
\newcommand{\RR}{\mathbb{R}}
\newcommand{\ZZ}{\mathbb{Z}}
\newcommand{\QQ}{\mathbb{Q}}
\newcommand{\Ga}{\Gamma}
\newcommand{\Bfun}{\mathrm{B}}
\newcommand{\im}{\operatorname{Im}}
\newcommand{\re}{\operatorname{Re}}
\newcommand{\lcm}{\operatorname{lcm}}
\newcommand{\SNF}{\operatorname{SNF}}
\newcommand{\Jac}{\operatorname{Jac}}
\newcommand{\rank}{\operatorname{rank}}
\newcommand{\Ind}{\operatorname{Index}}
\newcommand{\disc}{\operatorname{disc}}
\begin{document}
\begin{center}
{\LARGE\bfseries Theorem 2: Gamma-Normalization and Fractional Pi}\\[0.4em]
{\large algebraicity of normalized periods and the radical constants}\\[0.6em]
{\normalsize Isaev Iskhak Khamzatovich}\\[0.2em]
{\normalsize The <<Dynamic Principle>> Program --- the \texttt{hodge-laboratory}}\\[0.15em]
{\small Standalone theorem monograph · Монография 5/16}
\end{center}
\vspace{0.4em}
{\color{linkblue}\hrule height 0.8pt}
\vspace{0.8em}

\section{Setting}

The $\Ga$-normalization absorbs the transcendental factor of every
period, turning normalized periods into algebraic numbers of the
phase lattice $\frac1N\ZZ[\zeta_N]$. The program's requirement ---
fractional shares of $\pi$ instead of $4\pi^2$-normalizations ---
holds structurally on the stand.

\begin{definition}[$\Ga$-normalization]
For $P\in\frac{\Omega_{a,b}}{N}\ZZ[\zeta_N]$ the normalized period
is $\widehat{P}=P/\Omega_{a,b}\in\frac1N\ZZ[\zeta_N]$.
\end{definition}

\begin{theorem}[$\Ga$-normalization and fractional $\pi$]\label{th:norm}
Let $N\in\{15,30\}$. Then: (1) algebraicity:
$\widehat P\in\frac1N\ZZ[\zeta_N]$ --- the phase lattice has
denominator $N$; (2) reduction: each $\Omega_{a,b}$ is expressed
through canonical $\Ga$-values and factors $\pi/\sin(\pi k/N)$; at
$a+b=N$ --- a pure share of $\pi$; (3) no $4\pi^2$: $\pi$ enters
only via $\pi/\sin(\pi k/N)$ and the multiplication formula; no
power $4\pi^2$ arises in the normalization; (4) radicals:
\begin{align}
  b_{Ch}(15)&=1-\frac{1+\sqrt5+\sqrt3\,\sqrt{10-2\sqrt5}}{8},\\
  b_{Ch}(30)&=1-\frac{\sqrt{30+6\sqrt5}+\sqrt5-1}{8}.
\end{align}
\end{theorem}

\begin{proof}
(1) --- dividing the lattice $\frac{\Omega}N\ZZ[\zeta_N]$ by
$\Omega$. (2)--(3) --- the inventory: reflection introduces $\pi$
only as $\pi/\sin(\pi k/N)$, multiplication only as
$(2\pi)^{(m-1)/2}$ at integer $m$; the normalization is built from
the periods themselves, so $4\pi^2$ is not generated. (4):
$\cos24^\circ=\cos(60^\circ-36^\circ)$ with
$\cos36^\circ=\frac{1+\sqrt5}4$,
$\sin36^\circ=\frac{\sqrt{10-2\sqrt5}}4$ gives the level-15
radical; $\cos12^\circ=\cos(30^\circ-18^\circ)$ with
$\cos18^\circ=\frac{\sqrt{10+2\sqrt5}}4$,
$\sin18^\circ=\frac{\sqrt5-1}4$ gives the level-30 one.
Numerically: $0.086454$ and $0.021854$.
\end{proof}

\begin{remark}
The share of $\pi$ is always rational with denominator $N$: the
factor $\pi/\sin(\pi k/N)$ carries the angle $k\pi/N$; the phase
lattice carries the denominator $N$ in algebra. A structural claim
--- true for all blocks at once. The ladder $\pi/7$, $\pi/15$,
$\pi/30$ is one $\mu_N$-quantum ladder (certificates G, H).
\end{remark}

\begin{databox}[numbers]
$b_{Ch}(7)\approx0.7526$; $b_{Ch}(9)\approx0.1172$;
$b_{Ch}(11)\approx0.0823$; $b_{Ch}(15)\approx0.086454$;
$b_{Ch}(30)\approx0.021854$. The universal Tate twist $4\pi^2$ lives
in the lattice spectra (G1), not in the period normalization.
\end{databox}

\begin{statusbox}[summary]
The theorem fixes: the phase-lattice denominator equals the level
$N$; the transcendental part reduces to fractional $\pi$; the
radical constants are written in closed form.
\end{statusbox}

\begin{verifybox}[reproduction]
\texttt{python3 laboratory.py} $\to$ ``Stand N=15/30''; Lean:
\texttt{bch\_radicals}; the plot \texttt{fig\_bch.png}.
\end{verifybox}

\vfill
{\color{linkblue}\hrule height 0.6pt}
\vspace{0.5em}
{\small\color{gray}
License: individual exclusive license of Isaev Iskhak Khamzatovich.
All rights to the computations, formulas, and texts belong to the author.
Verification: \texttt{python3 laboratory.py} (``Stands''/``Certificates''),
Lean: \texttt{verification/lean/HodgeLaboratory.lean}.
}
\end{document}
